\subsection{Jump cut a předčasné ukončení skoku}

Při puštění tlačítka skoku během vzestupu se aktivuje jump cut, který předčasně ukončí skok a umožní hráči rychleji spadnout. Toto je důležitá mechanika, která umožňuje hráči kontrolovat výšku skoku podle délky stisknutí tlačítka.

Implementace využívá metodu \texttt{OnJumpRelease()}, která se registruje jako callback při uvolnění vstupu skoku. V této metodě se kontroluje, zda je hráčova vertikální rychlost kladná (tedy stoupá). Pokud ano, vynásobí se rychlost hodnotou \texttt{jumpCutVelocityMultiplier} z \texttt{PlayerControllerData}. Typicky se používá hodnota 0.5, což znamená, že při puštění tlačítka se rychlost sníží na polovinu.

V metodě \texttt{ApplyCustomGravity()} se pak při aktivním \texttt{jumpCut} aplikuje ještě vyšší gravitace než při normálním pádu. Toto zajistí okamžité ukončení vzestupu a rychlý přechod do pádu, což přispívá k responsivitě ovládání.

\begin{lstlisting}[language=csharp,caption={Listing 8: Metoda OnJumpRelease pro jump cut}]
private void OnJumpRelease(InputAction.CallbackContext ctx)
{
    if (rb != null && rb.linearVelocity.y > 0)
    {
        jumpCut = true;

        rb.linearVelocity = new Vector2(
            rb.linearVelocity.x,
            rb.linearVelocity.y * data.jumpCutVelocityMultiplier
        );
    }
}
\end{lstlisting}
